\section{An Enemy Made to Serve Life}

Augustine asks why the baptized still die when their sins have been forgiven. The question bears directly on Christian hope. Baptism is real regeneration; death is still before the regenerate. The believer receives a new life without receiving immediate exemption from the body's mortality. In \emph{The City of God} XIII.4--5, Augustine explains this remaining passage through the testing of faith and the extraordinary work of grace. Faith waits for what God has promised and what is not yet seen. Meanwhile, the punishment of sin is turned to the service of righteousness.

His sharpest example is martyrdom. The first disobedience brought death into human life. The martyr is threatened with death unless he disobeys, and accepts dying in order to remain faithful. The enemy is made to serve the obedience it once seemed able to destroy. Augustine is exact about the distinction: death has not become good in itself. God has granted the faithful a good use of what is evil. The wicked can misuse a good thing, as when the law becomes an occasion for transgression; the righteous, helped by grace, can use even death well.

Augustine's distinction explains how a martyr's death can be holy while the killing remains evil. The persecutor threatens the loss of bodily life to compel apostasy. By accepting death rather than denying Christ, the martyr keeps faith with the Lord who gave him that life and promises to raise his body. His constancy is sustained by grace. We honor the martyr because of his love and obedience, while condemning the violence inflicted on him.

Most Christians will not die by a persecutor's sentence. Yet Augustine's account helps us understand how death can become an obedient offering in an ordinary sickroom. The grace of Christ may meet us where independence gives way to dependence, familiar work becomes impossible, or a treasured plan must be relinquished. None of these losses becomes delightful merely because it can be offered. Their sorrow may remain. But sorrow need not determine the whole meaning of what we do with them. We can receive help, forgive, thank God for another's kindness, and consent to place unfinished work in his hands.

These acts are sustained by the same Lord who strengthens the martyrs. A person who has always managed the household may find being washed or fed more humiliating than an outsider can understand. A person who has found purpose in competent work may find enforced inactivity a severe trial. A parent may be troubled less by physical pain than by the prospect of leaving a child. We can bring the particular loss before God, ask for the help it requires, and seek to continue loving within it. The offering may need to be renewed each day as the same difficulty returns.

Approaching death, we may become unable to pray or act in the ways we had hoped. We can ask for the grace to accept help in that weakness. The person who can say little, follow only part of a familiar prayer, or simply consent to be prayed for remains within the Church's charity. Dependence on God has always been the truth of that life, including its strongest and most accomplished years. The help now received from others allows the weakened person to participate as he can, while the Church continues praying on his behalf.

We can practice this surrender before illness requires it. In confession we acknowledge sins we have tried to justify. In restitution we give up something we have wrongly treated as ours. In forgiveness we renounce revenge. In willing service we allow another person's need to interrupt our arrangements. These acts cannot guarantee a peaceful death; they help us become accustomed to obeying God when obedience costs us something. We will still need his grace when a loss is imposed on us and we have less freedom to choose its circumstances.

The baptized still die because the bodily fulfillment of Christ's promise remains future. Augustine explains that faith trusts God for this unseen reward. Death can therefore test that trust without defeating the grace received in Baptism. We rely on Christ to preserve us through the trial and to raise the body whose mortality he has allowed to remain.

